\chapter{INTRODUCTION}
\label{chap:intro}
\section{Overview}
Sequence models must frequently maintain a running summary of everything they have read: the current value of a variable in a program, the position of an object after a series of moves, or the state of a finite automaton after a string of inputs. This capability is known as \emph{state tracking}. Transformers address sequence modelling by attending over the full history at every step, which gives them great flexibility but makes their cost grow quadratically with sequence length. Linear recurrent neural networks, including linear attention \cite{lineartransformers} and selective state-space models \cite{mamba}, instead compress the history into a fixed-size state that is updated once per token. Their cost is linear in the sequence length and their memory during generation is constant, which makes them attractive for long-context and streaming applications.

Whether such a model can track state depends almost entirely on its \emph{state-transition matrix}, the matrix that maps the state at step $t-1$ to the state at step $t$. Early linear RNNs used diagonal transitions with non-negative entries, and a sequence of theoretical results has shown that this choice places hard limits on what they can track \cite{illusionofstate, sarrof, diagonalssm}. Two developments relaxed those limits. The delta rule \cite{fastweight, deltanet, gateddeltanet} replaces the diagonal transition with a generalized Householder matrix, a rank-one correction of the identity; allowing that correction to reach a reflection, which introduces negative eigenvalues, lets the model represent operations such as parity \cite{negeig}. DeltaProduct \cite{deltaproduct} then builds each transition as the product of $\nh$ generalized Householder factors, obtained by applying $\nh$ delta-rule updates per token, and shows that a larger $\nh$ lets the model track progressively harder groups, up to the symmetric group $S_5$.

In DeltaProduct the order $\nh$ is a single hyperparameter chosen by hand and applied identically to every token. Each additional factor adds computation for every token in the sequence. Yet different inputs plainly require different amounts of expressivity: a token that swaps two items is a simpler operation than one that cycles all five. A fixed global order must therefore be set for the hardest token the model will ever see, and that cost is paid on every token regardless.

This project replaces the fixed global order with a per-token order that the model learns. A small, monotone halting gate reads each token's representation and decides how many of the available factors to apply to it, and a compute-budget penalty encourages the gate to use as few as the task allows. Building this mechanism required answering a more basic question first: \emph{how many factors does a token actually require?} The answer depends not on the token alone but on the whole set of inputs the model is trained on, and that analysis forms a central part of this report.

\section{Motivation}
The practical motivation is efficiency. Every Householder factor in DeltaProduct costs an additional delta-rule update per token. When most tokens in a sequence require few factors and only some require many, a fixed order spends most of its computation on tokens that do not need it. An allocation that follows each token's requirement would pay the average requirement rather than the worst case.

The second motivation is that the correct value of $\nh$ is difficult to choose in advance. Chapter~\ref{chap:method} shows that the number of factors a token requires is determined by the algebraic structure of the entire input alphabet, and that adding a single, individually cheap input symbol can double the requirement of other symbols. A practitioner therefore cannot read the right order off the inputs one at a time. Choosing too small an order causes the model to fail on the tokens that exceed it; choosing too large an order wastes computation everywhere. A layer that learns its own order per token removes this hyperparameter from the design.

The third motivation is scientific. Adaptive computation has been studied extensively for the depth of Transformer and recurrent networks \cite{act, universaltransformer, pondernet, mod}, but not for the internal structure of the state transition of a linear RNN. Working at this level exposes a precise, verifiable relationship between the operation a token performs and the computation it needs.

\section{Organization of the Report}
Chapter~\ref{chap:lit} surveys the literature on linear recurrent architectures, the theory of state tracking and adaptive computation, and states the problem and objectives. Chapter~\ref{chap:method} formulates the task, derives the minimal order a token requires and describes the adaptive-order layer. Chapter~\ref{chap:results} presents the experimental setup and results, from the reproduction of the base method to the end-to-end training of the halting gate. Chapter~\ref{chap:conclusion} concludes and outlines future work. Appendix~I collects supplementary figures and tables, and the report closes with the key references and the project timeline.
